\documentclass[11pt,a4paper]{article}

\usepackage[spanish,es-noquoting]{babel}
\usepackage[utf8]{inputenc}
\usepackage[T1]{fontenc}
\usepackage{lmodern}
\usepackage{geometry}
\usepackage{booktabs}
\usepackage{tabularx}
\usepackage{longtable}
\usepackage{array}
\usepackage{multirow}
\usepackage{float}
\usepackage{graphicx}
\usepackage{enumitem}
\usepackage{hyperref}
\usepackage{xcolor}
\usepackage{colortbl}
\usepackage{fancyhdr}
\usepackage{amsmath}

\geometry{margin=2.5cm}
\setlength{\parskip}{0.5em}
\setlength{\parindent}{0pt}
\hypersetup{colorlinks=true, linkcolor=black, urlcolor=blue}

\definecolor{malignored}{RGB}{231,76,60}
\definecolor{benigngreen}{RGB}{46,204,113}
\definecolor{lightgray}{RGB}{245,245,245}
\definecolor{bestrow}{RGB}{212,239,223}
\definecolor{warnrow}{RGB}{253,237,236}
\definecolor{neutralrow}{RGB}{235,245,251}

\setlength{\headheight}{14pt}
\pagestyle{fancy}
\fancyhf{}
\lhead{\textit{Proyecto 2 -- Data Mining}}
\rhead{\textit{Semana 3: Validación, Optimización y Selección Final}}
\cfoot{\thepage}
\renewcommand{\headrulewidth}{0.4pt}

\title{\textbf{Proyecto 2 -- Data Mining}\\[0.4em]
       \large Semana 3: Validación, Optimización y Selección Final de Modelo}
\author{Javier España \and Ángel Esquit \and Roberto Barreda \and Diego López}
\date{Abril 2026}

\begin{document}

\maketitle
\thispagestyle{empty}

\begin{center}
\large Universidad del Valle de Guatemala\\[0.2em]
\normalsize Curso: Data Mining \quad|\quad Dataset: Wisconsin Breast Cancer Dataset\\[0.2em]
Variable objetivo: \texttt{Class} \;(2 = Benigno, 4 = Maligno)
\end{center}

\vspace{0.5em}\hrule\vspace{0.5em}

\tableofcontents
\newpage

% ─────────────────────────────────────────────
\section{Objetivo de Semana 3}
% ─────────────────────────────────────────────

Esta semana cerró el ciclo experimental del proyecto con cuatro frentes: (1) validación cruzada rigurosa y comparable de los tres modelos base; (2) optimización de hiperparámetros mediante \texttt{GridSearchCV} con la métrica clínica correcta; (3) diagnóstico de ajuste (sobreajuste vs.\ subajuste); y (4) selección formal del modelo final con justificación técnica explícita y cuantificación de trade-offs.

La métrica principal de toda esta semana fue el \textbf{Recall de la clase maligna} (\texttt{Class=4}), establecida desde el EDA de Semana 1 por el costo clínico de los falsos negativos.

% ─────────────────────────────────────────────
\section{Diseño Experimental}
% ─────────────────────────────────────────────

\subsection{Dataset y split}

Se trabajó con \texttt{Datos\_modelado.csv} (675 registros, 9 variables predictoras + \texttt{Class}), heredado de Semana 2 sin modificaciones para preservar comparabilidad. El split 80/20 estratificado con \texttt{random\_state=42} se mantuvo idéntico.

\begin{table}[H]
\centering
\caption{Verificación de proporciones de clase a lo largo del experimento}
\begin{tabular}{lccc}
\toprule
\textbf{Clase} & \textbf{Total} & \textbf{Train} & \textbf{Test} \\
\midrule
Benigna (2) & 65.04\% & 65.00\% & 65.19\% \\
Maligna (4) & 34.96\% & 35.00\% & 34.81\% \\
\bottomrule
\end{tabular}
\end{table}

La diferencia máxima entre particiones es 0.19 pp, confirmando que la estratificación no introduce sesgo de distribución.

\subsection{Esquema de validación cruzada}

Se adoptó \textbf{5-Fold Stratified Cross-Validation} (\texttt{StratifiedKFold(n\_splits=5, shuffle=True, random\_state=42)}) como método de validación. Este esquema garantiza:

\begin{itemize}[leftmargin=1.4cm]
\item Proporción de clases preservada en cada fold (crítico dado el desbalance 65/35).
\item Cinco estimaciones independientes del rendimiento, reduciendo la varianza del estimador.
\item Comparabilidad entre modelos: todos usan exactamente el mismo protocolo de folds.
\end{itemize}

\subsection{Control de fuga de información}

Todas las transformaciones (escalado con \texttt{StandardScaler}) se encapsulan dentro de un \texttt{sklearn.Pipeline}. Esto garantiza que el scaler se ajusta únicamente con los datos de entrenamiento de cada fold y nunca ve los datos de validación o de prueba, eliminando el riesgo de \textit{data leakage}.

% ─────────────────────────────────────────────
\section{Validación Base Pre-Tuning}
% ─────────────────────────────────────────────

\subsection{Resultados de validación cruzada -- Escenario A}

\begin{table}[H]
\centering
\caption{Cross-validation base en Escenario A (media $\pm$ desviación estándar, 5 folds)}
\small
\renewcommand{\arraystretch}{1.15}
\begin{tabular}{lcccc}
\toprule
\textbf{Modelo} & \textbf{Accuracy} & \textbf{Recall macro} & \textbf{F1} & \textbf{AUC} \\
\midrule
\rowcolor{bestrow}
Regresión Logística & $0.9704 \pm 0.0283$ & $0.9664 \pm 0.0351$ & $0.9702 \pm 0.0287$ & $0.9941 \pm 0.0074$ \\
KNN (k=5) & $0.9685 \pm 0.0224$ & $0.9636 \pm 0.0272$ & $0.9684 \pm 0.0226$ & $0.9843 \pm 0.0262$ \\
SVM (RBF) & $0.9611 \pm 0.0312$ & $0.9617 \pm 0.0341$ & $0.9613 \pm 0.0312$ & $0.9898 \pm 0.0109$ \\
\bottomrule
\end{tabular}
\normalsize
\end{table}

\subsection{Resultados de validación cruzada -- Escenario B}

\begin{table}[H]
\centering
\caption{Cross-validation base en Escenario B (top 3 variables)}
\small
\renewcommand{\arraystretch}{1.15}
\begin{tabular}{lcccc}
\toprule
\textbf{Modelo} & \textbf{Accuracy} & \textbf{Recall macro} & \textbf{F1} & \textbf{AUC} \\
\midrule
Regresión Logística & $0.9648 \pm 0.0271$ & $0.9583 \pm 0.0305$ & $0.9647 \pm 0.0272$ & $0.9912 \pm 0.0112$ \\
\rowcolor{bestrow}
KNN (k=5) & $0.9667 \pm 0.0239$ & $0.9634 \pm 0.0246$ & $0.9667 \pm 0.0237$ & $0.9823 \pm 0.0167$ \\
SVM (RBF) & $0.9574 \pm 0.0319$ & $0.9538 \pm 0.0339$ & $0.9575 \pm 0.0317$ & $0.9880 \pm 0.0122$ \\
\bottomrule
\end{tabular}
\normalsize
\end{table}

\subsection{Consistencia CV vs Test (base)}

\begin{table}[H]
\centering
\caption{Coherencia entre validación cruzada y conjunto de prueba -- modelos base}
\small
\renewcommand{\arraystretch}{1.15}
\begin{tabular}{llcccc}
\toprule
\textbf{Modelo} & \textbf{Esc.} & \textbf{Acc CV} & \textbf{Acc Test} & \textbf{Gap Acc} & \textbf{Gap Recall} \\
\midrule
Regresión Logística & A & 0.9704 & 0.9556 & $-$0.0148 & $-$0.0153 \\
KNN & A & 0.9685 & 0.9556 & $-$0.0130 & $-$0.0126 \\
SVM & A & 0.9611 & 0.9630 & $+$0.0019 & $\approx 0$ \\
\midrule
\rowcolor{warnrow}
Regresión Logística & B & 0.9648 & 0.9259 & $-$0.0389 & $-$0.0498 \\
KNN & B & 0.9667 & 0.9481 & $-$0.0185 & $-$0.0230 \\
SVM & B & 0.9574 & 0.9481 & $-$0.0093 & $-$0.0134 \\
\bottomrule
\end{tabular}
\normalsize
\end{table}

El gap más grande lo muestra LR en Escenario B ($-$0.0498 en Recall), consistente con la caída observada en Semana 2. En Escenario A todos los gaps son menores a 0.016, lo que indica generalización estable y habilita una etapa de tuning confiable.

% ─────────────────────────────────────────────
\section{Optimización de Hiperparámetros: LR y KNN}
% ─────────────────────────────────────────────

\subsection{Metodología}

Se ejecutó \texttt{GridSearchCV} con el mismo \texttt{StratifiedKFold(5)} y como métrica de selección el \textbf{Recall de la clase maligna} (\texttt{pos\_label=4}). El test set no participó en ninguna decisión de tuning.

\subsection{Mejores hiperparámetros encontrados}

\begin{table}[H]
\centering
\caption{Mejores hiperparámetros de LR y KNN por escenario}
\renewcommand{\arraystretch}{1.2}
\begin{tabularx}{\textwidth}{l l c c X}
\toprule
\textbf{Esc.} & \textbf{Modelo} & \textbf{Best Recall CV} & \textbf{Std} & \textbf{Mejores parámetros} \\
\midrule
A & LR & 0.9789 & 0.0421 & \texttt{C=0.1, penalty=l1, solver=liblinear, class\_weight=balanced} \\
A & KNN & 0.9632 & 0.0488 & \texttt{n\_neighbors=9, weights=uniform, metric=euclidean} \\
\midrule
B & LR & 0.9471 & 0.0333 & \texttt{C=3, penalty=l2, solver=lbfgs, class\_weight=balanced} \\
B & KNN & 0.9578 & 0.0268 & \texttt{n\_neighbors=7, weights=uniform, metric=euclidean} \\
\bottomrule
\end{tabularx}
\end{table}

\textbf{Observaciones sobre LR:} el mejor $C$ en Escenario A es 0.1 (alta regularización), lo que confirma que el modelo se beneficia de penalización fuerte cuando trabaja con variables colineales. En Escenario B, el mejor $C$ es 3 (menor regularización), pues con solo 3 variables hay menos redundancia que penalizar.

\subsection{Comparación base vs optimizado: LR y KNN}

\begin{table}[H]
\centering
\caption{Impacto del tuning en Recall de maligna y FN -- LR y KNN}
\renewcommand{\arraystretch}{1.15}
\begin{tabular}{llccccc}
\toprule
\textbf{Esc.} & \textbf{Modelo} & \textbf{Tipo} & \textbf{Recall CV} & \textbf{Recall Test} & \textbf{FN test} & \textbf{$\Delta$ FN} \\
\midrule
\multirow{4}{*}{A}
 & LR & Base & 0.9526 & 0.9362 & 3 & \multirow{2}{*}{$-1$} \\
 & \rowcolor{bestrow}LR & Optimizado & 0.9789 & 0.9574 & 2 & \\
 & KNN & Base & 0.9472 & 0.9362 & 3 & \multirow{2}{*}{$+1$} \\
 & \rowcolor{warnrow}KNN & Optimizado & 0.9632 & 0.9149 & 4 & \\
\midrule
\multirow{4}{*}{B}
 & LR & Base & 0.9366 & 0.8511 & 7 & \multirow{2}{*}{$-3$} \\
 & \rowcolor{bestrow}LR & Optimizado & 0.9471 & 0.9149 & 4 & \\
 & KNN & Base & 0.9525 & 0.9149 & 4 & \multirow{2}{*}{$0$} \\
 & KNN & Optimizado & 0.9578 & 0.9149 & 4 & \\
\bottomrule
\end{tabular}
\end{table}

\textbf{Hallazgo crítico:} el tuning de KNN en Escenario A \textit{empeora} el recall en test ($-$0.0213) y aumenta los FN de 3 a 4. Esto indica que los hiperparámetros encontrados en validación no se transfieren bien a test para KNN, lo que descalifica a KNN como candidato final.

\subsection{Diagnóstico de ajuste -- LR y KNN}

\begin{table}[H]
\centering
\caption{Análisis de sobreajuste/subajuste -- modelos optimizados LR y KNN}
\renewcommand{\arraystretch}{1.15}
\begin{tabular}{llcccc}
\toprule
\textbf{Esc.} & \textbf{Modelo} & \textbf{Recall train} & \textbf{Recall valid} & \textbf{Gap} & \textbf{Diagnóstico} \\
\midrule
A & LR & 0.9841 & 0.9789 & 0.0052 & Estable \\
A & KNN & 0.9643 & 0.9632 & 0.0011 & Estable \\
B & LR & 0.9471 & 0.9471 & 0.0000 & Estable \\
B & KNN & 0.9630 & 0.9578 & 0.0052 & Estable \\
\bottomrule
\end{tabular}
\end{table}

Los gaps train-valid son mínimos en todos los modelos. Sin embargo, KNN en Escenario A exhibe una discordancia entre la estabilidad en CV y la caída en test, lo que sugiere sensibilidad a la distribución específica del fold o del test set.

% ─────────────────────────────────────────────
\section{Optimización de Hiperparámetros: SVM}
% ─────────────────────────────────────────────

\subsection{Espacio de búsqueda}

\begin{table}[H]
\centering
\caption{Espacio de búsqueda del GridSearchCV para SVM}
\begin{tabular}{lll}
\toprule
\textbf{Hiperparámetro} & \textbf{Valores explorados} & \textbf{Descripción} \\
\midrule
\texttt{C} & 0.1, 0.3, 1, 3, 10, 30, 100 & Control de margen vs error de clasificación \\
\texttt{gamma} & \texttt{`scale'}, 0.01, 0.03, 0.1, 0.3, 1 & Alcance del kernel RBF \\
\texttt{class\_weight} & \texttt{None}, \texttt{`balanced'} & Compensación de desbalance de clases \\
\bottomrule
\end{tabular}
\end{table}

\subsection{Mejores configuraciones SVM}

\begin{table}[H]
\centering
\caption{Mejores configuraciones SVM por escenario (top resultados del GridSearchCV)}
\renewcommand{\arraystretch}{1.1}
\begin{tabular}{lcccccc}
\toprule
\textbf{Esc.} & \textbf{C} & \textbf{Gamma} & \textbf{class\_weight} & \textbf{Recall CV} & \textbf{Std CV} & \textbf{Rank} \\
\midrule
A & 0.1 & 1 & None & 1.0000 & 0.0000 & 1 \\
A & 0.1 & 1 & balanced & 1.0000 & 0.0000 & 1 \\
A & 0.3 & 1 & None & 0.9947 & 0.0105 & 3 \\
\midrule
B & 3 & scale & balanced & 0.9735 & 0.0166 & 1 \\
B & 10 & scale & balanced & 0.9735 & 0.0236 & 1 \\
B & 10 & 0.3 & balanced & 0.9735 & 0.0236 & 1 \\
\bottomrule
\end{tabular}
\end{table}

\subsection{Resultado central: base vs optimizado}

\begin{table}[H]
\centering
\caption{SVM base vs optimizado -- comparación completa de métricas en test}
\renewcommand{\arraystretch}{1.2}
\begin{tabular}{llccccccc}
\toprule
\textbf{Esc.} & \textbf{Tipo} & \textbf{Recall CV} & \textbf{Recall Test} & \textbf{Prec.\ Test} & \textbf{F1 Test} & \textbf{AUC Test} & \textbf{FN} & \textbf{TN/FP/FN/TP} \\
\midrule
A & Base & 0.9632 & 0.9574 & 0.9375 & 0.9474 & 0.9908 & 2 & 85/3/2/45 \\
\rowcolor{bestrow}
A & Optimizado & \textbf{1.0000} & \textbf{1.0000} & 0.9038 & 0.9495 & 0.9744 & \textbf{0} & 83/5/0/47 \\
\midrule
B & Base & 0.9418 & 0.9149 & 0.9348 & 0.9247 & 0.9877 & 4 & 85/3/4/43 \\
\rowcolor{neutralrow}
B & Optimizado & 0.9735 & 0.9787 & 0.9388 & 0.9583 & 0.9705 & 1 & 85/3/1/46 \\
\bottomrule
\end{tabular}
\end{table}

\textbf{Resultado excepcional en Escenario A:} el SVM optimizado logra Recall = 1.0000 y FN = 0 en el conjunto de prueba. Esto significa que \textbf{ningún caso maligno fue clasificado incorrectamente}. El costo de este logro es una caída en precisión de 0.9375 a 0.9038 (5 falsos positivos en lugar de 3) y en AUC de 0.9908 a 0.9744.

\subsection{Diagnóstico de ajuste -- SVM}

\begin{table}[H]
\centering
\caption{Análisis de sobreajuste/subajuste -- SVM optimizado}
\renewcommand{\arraystretch}{1.15}
\begin{tabular}{lcccccc}
\toprule
\textbf{Esc.} & \textbf{Tipo} & \textbf{Recall train} & \textbf{Recall valid} & \textbf{Gap} & \textbf{Diagnóstico} & \textbf{Sostiene en test} \\
\midrule
A & Base & 0.9749 & 0.9632 & 0.0117 & Estable & Referencia \\
\rowcolor{bestrow}
A & Optimizado & 1.0000 & 1.0000 & 0.0000 & Estable & Sí \\
B & Base & 0.9458 & 0.9418 & 0.0040 & Estable & Referencia \\
\rowcolor{neutralrow}
B & Optimizado & 0.9762 & 0.9735 & 0.0026 & Estable & Sí \\
\bottomrule
\end{tabular}
\end{table}

El gap train-valid de SVM A optimizado es exactamente 0.0000, lo que podría interpretarse como indicio de sobreajuste. Sin embargo, el hecho de que \textit{este mismo rendimiento se sostenga en test} descarta el sobreajuste clásico. La explicación es más matizada: el dataset presenta separabilidad alta y el kernel RBF con parámetros específicos ($C=0.1$, $\gamma=1$) encontró una frontera de decisión que cubre perfectamente todos los casos malignos. Este comportamiento es esperable en datasets bien condicionados con alta separación inter-clase.

% ─────────────────────────────────────────────
\section{Comparación Final: Todos los Modelos Optimizados}
% ─────────────────────────────────────────────

\begin{table}[H]
\centering
\caption{Ranking final de los seis candidatos optimizados (ordenado por Recall de maligna en test)}
\small
\renewcommand{\arraystretch}{1.2}
\begin{tabular}{llccccccc}
\toprule
\textbf{Esc.} & \textbf{Modelo} & \textbf{Recall CV} & \textbf{Recall Test} & \textbf{Prec.\ Test} & \textbf{F1 Test} & \textbf{AUC Test} & \textbf{FN} & \textbf{Sostiene} \\
\midrule
\rowcolor{bestrow}
A & \textbf{SVM} & 1.0000 & \textbf{1.0000} & 0.9038 & 0.9495 & 0.9744 & \textbf{0} & Sí \\
A & LR & 0.9789 & 0.9574 & 0.9375 & 0.9474 & \textbf{0.9944} & 2 & Sí \\
\rowcolor{neutralrow}
B & SVM & 0.9735 & 0.9787 & 0.9388 & \textbf{0.9583} & 0.9705 & 1 & Sí \\
A & KNN & 0.9632 & 0.9149 & 0.9348 & 0.9247 & 0.9889 & 4 & No \\
\rowcolor{warnrow}
B & LR & 0.9471 & 0.9149 & 0.9348 & 0.9247 & 0.9920 & 4 & No \\
\rowcolor{warnrow}
B & KNN & 0.9578 & 0.9149 & 0.9348 & 0.9247 & 0.9897 & 4 & No \\
\bottomrule
\end{tabular}
\normalsize
\end{table}

\textbf{Lectura del ranking:}
\begin{itemize}[leftmargin=1.4cm]
\item \textbf{SVM A optimizado} domina en el criterio clínico principal (Recall, FN). Es el único candidato que elimina completamente los falsos negativos.
\item \textbf{LR A optimizado} es el segundo mejor y el más fuerte en AUC (0.9944), lo que lo hace preferible si el contexto requiriera probabilidades bien calibradas.
\item \textbf{KNN} en ningún escenario supera su baseline de forma consistente en test; queda descartado como modelo final.
\item Los tres modelos del Escenario B que ``no sostienen'' en test (KNN-A, LR-B, KNN-B) exhiben una discordancia entre la mejora en CV y la no-mejora en test, señal de que el tuning no generalizó en esos casos.
\end{itemize}

% ─────────────────────────────────────────────
\section{Modelo Final Seleccionado}
% ─────────────────────────────────────────────

\subsection{Selección}

\textbf{Modelo final: SVM optimizado con kernel RBF, Escenario A (todas las variables).}

Parámetros finales: \texttt{C=0.1}, \texttt{gamma=1}, \texttt{kernel=`rbf'}, \texttt{class\_weight=None}, \texttt{StandardScaler}.

\subsection{Justificación técnica}

\begin{enumerate}[leftmargin=1.4cm]
\item \textbf{Recall test = 1.0000:} máximo posible. Ningún caso maligno clasificado incorrectamente.
\item \textbf{FN test = 0:} criterio clínico dominante del proyecto.
\item \textbf{Sostiene en test:} la mejora observada en CV se confirma en el conjunto de prueba independiente.
\item \textbf{Sin sobreajuste clásico:} el gap train-valid de 0.0 se explica por separabilidad del dataset, no por memorización.
\end{enumerate}

\subsection{Cuantificación del trade-off clínico}

\begin{table}[H]
\centering
\caption{Trade-off clínico: SVM A optimizado vs LR A optimizado (conjunto de prueba)}
\renewcommand{\arraystretch}{1.2}
\begin{tabular}{lcccc}
\toprule
\textbf{Modelo} & \textbf{TP (malignos detectados)} & \textbf{FN (omitidos)} & \textbf{FP (falsos alarmas)} & \textbf{AUC} \\
\midrule
\rowcolor{bestrow}
SVM A optimizado & 47 & \textbf{0} & 5 & 0.9744 \\
LR A optimizado & 45 & 2 & 3 & 0.9944 \\
\midrule
\textbf{Diferencia} & $+2$ & $-2$ & $+2$ & $-0.0200$ \\
\bottomrule
\end{tabular}
\end{table}

SVM A evita 2 falsos negativos a cambio de 2 falsos positivos adicionales y 0.02 de AUC. En el contexto clínico: un falso negativo puede significar un cáncer no tratado; un falso positivo implica una biopsia adicional o seguimiento. La relación costo-beneficio es claramente favorable a SVM A desde la perspectiva médica.

% ─────────────────────────────────────────────
\section{Sesgo, Varianza y Estabilidad del Modelo Final}
% ─────────────────────────────────────────────

El análisis de sesgo-varianza del SVM A optimizado muestra:

\begin{itemize}[leftmargin=1.4cm]
\item \textbf{Sesgo bajo:} el modelo clasifica correctamente la totalidad de los casos en validación (Recall CV = 1.0), indicando que tiene suficiente complejidad para aprender la estructura del problema.
\item \textbf{Varianza baja:} la desviación estándar del Recall en CV es 0.0, lo que significa que los cinco folds obtuvieron exactamente el mismo resultado. Esto es posible porque el dataset tiene separabilidad alta y los 540 registros de entrenamiento son suficientes para que el modelo sea estable.
\item \textbf{Generalización:} la transferencia perfecta de CV a test (Recall 1.0 en ambos) descarta memorización. Si hubiera sobreajuste, esperaríamos que el test degradara respecto al CV.
\end{itemize}

La estabilidad del modelo es también coherente con el contexto del dataset: alta separabilidad inter-clase, rangos acotados (1--10) y sin ruido severo hacen que la frontera de decisión sea robusta.

% ─────────────────────────────────────────────
\section{Verificación de Hipótesis de Semana 1}
% ─────────────────────────────────────────────

\begin{table}[H]
\centering
\caption{Estado de las hipótesis formuladas en el EDA al cierre de Semana 3}
\renewcommand{\arraystretch}{1.3}
\begin{tabularx}{\textwidth}{l l X}
\toprule
\textbf{Hipótesis} & \textbf{Veredicto} & \textbf{Evidencia} \\
\midrule
\rowcolor{bestrow}
H1: Bare Nuclei y variables de uniformidad serán las más importantes & Confirmada parcialmente & El Escenario B con esas 3 variables alcanza Recall 0.979; el modelo final usa todas las variables \\
H2: Modelos no lineales capturan mejor relaciones & Confirmada & SVM-RBF supera a LR en Recall en ambos escenarios tras tuning \\
\rowcolor{bestrow}
H3: Recall de maligna como métrica principal & Confirmada y aplicada & Fue la métrica objetivo de todo el GridSearchCV de Semana 3 \\
H4: Colinealidad puede desestabilizar LR & Confirmada & LR pierde 8.5 pp de Recall en Escenario B; KNN y SVM solo 4.3 pp o menos \\
\rowcolor{bestrow}
H5: No eliminar outliers clínicos & Sostenida & El modelo final entrenado con outliers preservados logra Recall = 1.0 \\
\bottomrule
\end{tabularx}
\end{table}

% ─────────────────────────────────────────────
\section{Conclusiones de Semana 3}
% ─────────────────────────────────────────────

\begin{enumerate}[leftmargin=1.4cm]
\item El diseño experimental (split heredado, CV estratificada, mismo protocolo para todos los modelos) fue estable y comparable de inicio a fin.
\item \textbf{SVM con kernel RBF optimizado en Escenario A} es el modelo final: Recall = 1.0, FN = 0, sin evidencia de sobreajuste.
\item LR optimizada queda como alternativa fuerte para contextos donde se requiera AUC alto o probabilidades calibradas.
\item KNN no muestra mejora consistente tras tuning y no es un candidato viable para producción.
\item Las 5 hipótesis del EDA quedaron verificadas o confirmadas, cerrando la coherencia del proyecto desde Semana 1.
\item El trade-off del modelo final (2 FP adicionales a cambio de 0 FN) es clínicamente defendible y explícitamente cuantificado.
\end{enumerate}

\vspace{1em}
\noindent\textbf{Repositorio del proyecto:}
\url{https://github.com/AngelEsquit/PRY2-MD.git}

\end{document}